\chapter{Conclusions}
 
This thesis proposed a \ac{DBN} framework for probabilistic short-horizon prediction of human driver maneuvers from window-based highway observations. 
The results show that the proposed model can capture meaningful temporal behavior patterns and provide useful maneuver forecasts in a causal online setting.

\par Among the compared configurations, the hierarchical emission model gave the best overall results and was selected as the final setup. 
It consistently outperformed the \ac{PoE} alternative and provided the best final balance between model fit, latent-state usage, and interpretability. 
The selected model achieved moderate performance in strict exact next-step prediction, which confirms that fine-grained maneuver discrimination is difficult. 
At the same time, the short-horizon results were clearly stronger, especially when semantically similar maneuvers were grouped. 
This shows that many errors are near-misses between related behaviors rather than completely wrong predictions. 
The horizon-based evaluation further supported this finding, with a strong hit rate over the 2\,s prediction horizon.

\par The latent-cardinality ablation study further supported the final design choice. 
Using only three action states made the behavioral partition too coarse, while removing the style variable caused a much larger drop in model quality. 
Although five action states improved the numerical fit, the resulting state space became less compact and harder to interpret. 
Overall, the selected two-style, four-action structure provided the best trade-off between prediction quality, interpretability, and model complexity.

\par Finally, the runtime benchmark showed that the framework is suitable for online use. 
The measured end-to-end latency remained far below the available decision interval, which makes the proposed approach practical for real-time short-horizon maneuver prediction.

\section{Future Work}

Future work should focus on four main directions.  
First, the observation design can be improved by investigating alternative feature sets and grouping strategies in order to achieve cleaner separation between latent states. 
Second, richer driving datasets with more diverse maneuvers should be considered, since the selected highD dataset mainly contains lane-following highway behavior. 
Third, alternative latent-state structures could be explored to represent driving style and maneuver dynamics more effectively. 
Finally, the proposed predictor can be integrated into a larger \ac{AV} system, where its probabilistic maneuver forecasts may support planning and risk-assessment modules.